\textLigne{Entraînement MathALEA -- Généralités}
\medskip

\qrwithlabel{Lire une image}{https://coopmaths.fr/alea/?uuid=6c6b3&id=1F1-2&uuid=b8946&id=1F1-11&lang=fr-CH&v=eleve&es=1011010&title=Lire+une+image+\%C3\%A0+partir+du+graphique}
\qrwithlabel{Utiliser le vocabulaire lié aux fonctions}{https://coopmaths.fr/alea/?uuid=0eecd&id=1F1-5&n=3&d=10&s=3&cd=0&uuid=4daef&id=1F1-6&n=5&d=10&s=5&cd=1&uuid=b92da&id=1F1-4&lang=fr-CH&v=eleve&es=1011010&title=Vocabulaire+des+fonctions}
\qrwithlabel{Calculer une image ou une préimage}{https://coopmaths.fr/alea/?uuid=c9382&id=1F1-7&n=3&d=10&s=3&s2=1&s3=5&cd=1&uuid=8a78e&id=1F1-8&n=2&d=10&s=1&s2=1&s3=4&cd=1&uuid=ba520&id=1F1-3&n=3&d=10&s=1-2&s2=2&s3=5&cd=1&lang=fr-CH&v=eleve&es=1011010&title=Calculs+d\%27images+et+de+pr\%C3\%A9images}
\qrwithlabel{Compléter un tableau de valeurs}{https://coopmaths.fr/alea/?uuid=afb2f&id=1F1-10&n=3&d=10&s=5&cd=1&lang=fr-CH&v=eleve&es=0011000&title=Compl\%C3\%A9ter+un+tableau+de+valeurs}

\medskip

\textLigne{Entraînement MathALEA -- Fonctions affines et droite}
\medskip

\qrwithlabel{Déterminer la pente de la droite passant par deux points}{https://coopmaths.fr/alea/?uuid=1ea16&id=1F2-1&lang=fr-CH&v=eleve&es=0011010&title=D\%C3\%A9terminer+la+pente+d\%27une+droite+passant+par+deux+points}
\qrwithlabel{Déterminer l'équation d'une droite représentée graphiquement}{https://coopmaths.fr/alea/?uuid=41e6f&id=1F2-6&n=3&d=10&s=2&cd=1&lang=fr-CH&v=eleve&es=0011010&title=D\%C3\%A9terminer+l\%27\%C3\%A9quation+d\%27une+droite+repr\%C3\%A9sent\%C3\%A9e+graphiquement}
\qrwithlabel{Déterminer l'équation d'une droite connaissant sa pente et un point}{https://coopmaths.fr/alea/?uuid=d1da3&id=1F2-5&lang=fr-CH&v=eleve&es=0011000&title=D\%C3\%A9terminer+l\%27\%C3\%A9quation+connaissant+la+pente+et+un+point}
\qrwithlabel{Déterminer l'équation d'une droite passant par deux points}{https://coopmaths.fr/alea/?uuid=0cee9&id=1F2-2&uuid=1bb30&id=1F2-3&lang=fr-CH&v=eleve&es=1011010&title=D\%C3\%A9terminer+l\%27\%C3\%A9quation+d\%27une+droite+\%C3\%A0+partir+de+deux+points}
\medskip
\medskip

\qrwithlabel{Points alignés}{https://coopmaths.fr/alea/?uuid=b1777&id=1F2-7&n=3&d=10&s=1&cd=1&lang=fr-CH&v=eleve&es=0011000&title=Points+align\%C3\%A9s}
\qrwithlabel{Point d'intersection de deux droites}{https://coopmaths.fr/alea/?uuid=1b828&id=1F2-9&n=1&d=10&s=1&cd=1&uuid=4b211&id=1F2-11&uuid=8ff15&id=1F2-10&n=1&d=10&s=2&s2=1&cd=1&lang=fr-CH&v=eleve&es=1011000&title=Points+d\%27intersection}
\qrwithlabel{Droites parallèles et perpendiculaires}{https://coopmaths.fr/alea/?uuid=2db22&id=1F2-12&n=2&d=10&s=3&s2=3&cd=1&lang=fr-CH&v=eleve&es=0011000&title=Droites+parall\%C3\%A8les+et+perpendiculaires}
\medskip

\textLigne{Entraînement MathALEA -- Fonctions quadratiques et paraboles}
\medskip

\qrwithlabel{Lecture graphique}{https://coopmaths.fr/alea/?uuid=a896e&id=1F3-1&n=3&d=10&s=4&cd=1&lang=fr-CH&v=eleve&es=0011000&title=Lecture+graphique}
\qrwithlabel{Forme canonique}{https://coopmaths.fr/alea/?uuid=60504&id=1F3-2&n=2&d=10&cd=1&lang=fr-CH&v=eleve&es=0011000&title=Forme+canonique}
\qrwithlabel{Sens de variation}{https://coopmaths.fr/alea/?uuid=16f97&id=1F3-4&n=4&d=10&s=1-3-7-5&cd=1&lang=fr-CH&v=eleve&es=0011000&title=Sens+de+variation}
\qrwithlabel{Points d'intersection}{https://coopmaths.fr/alea/?uuid=e37e2&id=1F2-13&n=3&d=10&s=2&s2=false&s3=4&cd=1&uuid=e37e2&id=1F2-13&n=3&d=10&s=2&s2=true&s3=4&cd=1&lang=fr-CH&v=eleve&es=0211001&title=Points+d\%27intersection}
\medskip

\qrwithlabel{Retrouver l'équation d'une parabole}{https://coopmaths.fr/alea/?uuid=392b3&id=1F3-6&lang=fr-CH&v=eleve&es=0011000&title=\%C3\%89quation+de+la+parabole+connaissant+le+sommet+et+un+point}
\qrwithlabel{Étude complète d'une parabole}{https://coopmaths.fr/alea/?uuid=5ea23&id=1F3-3&lang=fr-CH&v=eleve&es=0011000&title=\%C3\%89tude+compl\%C3\%A8te+d\%27une+parabole}
